\documentclass[12pt,a4paper,oneside]{scrreprt}
\usepackage[utf8]{inputenc}
\usepackage[T1]{fontenc}
\usepackage[english,ngerman]{babel}
\usepackage{amsmath,amssymb,amsfonts}
\usepackage{booktabs}
\usepackage{array}
\usepackage{etoolbox}
\AtBeginEnvironment{tabular}{\small}   % keep tables within the text width
\usepackage{listings}
\lstdefinestyle{promptbox}{
    basicstyle=\footnotesize\ttfamily,
    frame=single,
    framesep=6pt,
    breaklines=true,
    columns=fullflexible,
    keepspaces=true,
    xleftmargin=1.5em,
    xrightmargin=1.5em,
    aboveskip=1.4em,
    belowskip=0.6em,
    captionpos=b
}
\usepackage{tikz}
\usetikzlibrary{arrows.meta,positioning,shapes.geometric,fit,backgrounds,calc}
\usepackage{graphicx}
\usepackage{geometry}
% Page layout tuned to roughly 1500 characters per full text page
\geometry{a4paper, top=2.5cm, bottom=3cm, left=4.5cm, right=4cm,
          headsep=0.8cm, footskip=1.2cm}
\usepackage{setspace}
\doublespacing

% --- Section numbering and table of contents depth ---
\setcounter{secnumdepth}{3}   % number down to subsubsections
\setcounter{tocdepth}{2}      % list chapters, sections, subsections in the TOC

% --- Heading layout: clear spacing so sections stand out ---
\RedeclareSectionCommand[beforeskip=-0.5\baselineskip, afterskip=2\baselineskip]{chapter}
\RedeclareSectionCommand[beforeskip=-2\baselineskip, afterskip=0.8\baselineskip]{section}
\RedeclareSectionCommand[beforeskip=-1.6\baselineskip, afterskip=0.6\baselineskip]{subsection}
\RedeclareSectionCommand[beforeskip=-1.2\baselineskip, afterskip=0.4\baselineskip]{subsubsection}
\setkomafont{chapter}{\normalfont\sffamily\bfseries\huge}
\setkomafont{section}{\normalfont\sffamily\bfseries\Large}
\setkomafont{subsection}{\normalfont\sffamily\bfseries\large}
\setkomafont{subsubsection}{\normalfont\sffamily\bfseries\normalsize}

% --- Running heads and page numbers (KOMA-native, replaces fancyhdr) ---
\usepackage[automark,headsepline]{scrlayer-scrpage}
\clearpairofpagestyles
\ohead{\headmark}
\ofoot*{\pagemark}   % page number bottom right on every page
\pagestyle{scrheadings}

\usepackage{acronym}
\usepackage{natbib}
\usepackage[hidelinks,bookmarksnumbered,bookmarksopen,bookmarksopenlevel=1]{hyperref}
\usepackage{bookmark}

\begin{document}

\pagenumbering{roman}

% Title Page
\begin{titlepage}
    \centering
    \vspace*{1cm}
    % \includegraphics[width=0.5\textwidth]{goethe_logo.png} % Placeholder for the Goethe University logo
    \vspace{1cm}
    {\scshape\LARGE Goethe-Universität Frankfurt am Main \par}
    \vspace{0.5cm}
    {\scshape\Large Fachbereich Informatik \par}
    \vspace{2cm}
    {\huge\bfseries Instruction Robustness in RAG/GraphRAG Agent Pipelines \par}
    \vspace{1.5cm}
    {\Large\itshape Bachelor Thesis \par}
    \vspace{2cm}
    \begin{flushleft}
        \large
        \textbf{Name:} Christian Block \\
        \textbf{Matriculation Number:} 6522508 \\
        \textbf{First examiner:} Prof. Dr. Visvanathan Ramesh \\
        \textbf{Date of Submission:}
    \end{flushleft}
    \vfill
\end{titlepage}

% Declaration
\selectlanguage{ngerman}
\chapter*{Erklärung zur Abschlussarbeit}
gemäß § 35, Abs. 16 der Ordnung für den Bachelorstudiengang Informatik vom 17. Juni 2019:\\
Hiermit erkläre ich (Nachname, Vorname) \textbf{Block, Christian}, dass ich die vorliegende Arbeit selbstständig und ohne Benutzung anderer als der angegebenen Quellen und Hilfsmittel verfasst habe. Ich bestätige außerdem, dass die vorliegende Arbeit nicht, auch nicht auszugsweise, für eine andere Prüfung oder Studienleistung verwendet wurde. Zudem versichere ich, dass alle eingereichten schriftlichen gebundenen Versionen meiner vorliegenden Bachelorarbeit mit der digital eingereichten elektronischen Version meiner Bachelorarbeit übereinstimmen.\\
Frankfurt am Main, den DATUM \hfill \rule{5cm}{0.4pt}\\
\hspace*{0pt}\hfill Unterschrift der/des Studierenden

\selectlanguage{english}
\chapter*{Abstract}

Retrieval-augmented generation (RAG) combines information retrieval with language generation, so that a language model can condition its output on external and updatable evidence. The system prompt specifies the instructions, constraints, and output
format the generator should adhere to. Because it is written in natural language, its effect on system behaviour is hard to predict and hard to verify. Prior work on prompt sensitivity in language models has studied that effect outside retrieval-augmented systems. This thesis addresses the resulting gap by asking to what extent systematic changes in the system prompt influence the behaviour and outcomes of a RAG system.

The thesis builds a formal model of a RAG system based on the formulation of \citet{lewis2020rag}. A retriever assigns probabilities
$p_\eta(z\mid q,D)$ to documents from the knowledge base, and  a generator produces the answer token by token. The system prompt $S$ enters as an explicit variable, $p_\theta(y_i\mid q,C,y_{1:i-1},S)$. It is modelled as a set of checkable
requirements $\Sigma$ written in a wording $\varphi$, such that $S=\rho(\Sigma,\varphi)$. Each of three published baseline prompts is written in four wordings: its canonical form and three rewordings, each of which changes either the vocabulary, the sentence structure or the layout. Every rewording leaves the requirements fixed and changes only $\varphi$.

The properties a
good answer must have are fixed beforehand: it must be grounded in the retrieved evidence, relevant to the question asked, correct with respect to the reference answer, and compliant with the instructions it was given. The knowledge base and the query set come from an established question-answer
dataset, so that the measurement is reproducible. Only the system prompt is manipulated. The language model, the
embedding model, the vector index, the chunking strategy and the retrieval
depth are held fixed. The generator therefore receives identical passages in every condition. Any difference in the output can be attributed to the generator alone. Generations are repeated to account for decoding
randomness.\\

The null hypothesis of
equal behaviour across wordings is rejected. The clearest effects are on
instruction adherence, which three rewordings raised by up to 26 percentage points and one
lowered by up to 12. Correctness, measured by ROUGE-L, BLEU-1 and BERTScore, changed by a few percentage points, mostly because the answers changed in length. Faithfulness and answer relevance showed no consistent change. 


\selectlanguage{english}
\chapter*{Acknowledgments}
I want to express my sincere gratitude to Professor Ramesh for his guidance and mentorship throughout this thesis and Professor Gao for being the second supervisor of this thesis.

\tableofcontents

\cleardoublepage
\pagenumbering{arabic}

\chapter{Introduction}

\section{Background and Motivation}

Retrieval-Augmented Generation (RAG) systems combine information retrieval with
language generation. A RAG
system takes the user query, retrieves relevant information from an external
knowledge base and gives that information to a
language model with instructions for generating the final response.

The system prompt is where the intended behaviour of the system is specified, and it shapes the output the system produces. Whereas the retrieved evidence varies with the query and the knowledge base, the system prompt specifies how the model should behave, which instructions and constraints
apply, and which output format it should produce. Because it is written in natural language, its effect on system behaviour is difficult to
predict. Its influence therefore has to be examined
through controlled experiments.

\section{Research Objective}
The central research question is the following: \\
\textit{To what extent do
systematic changes in the system prompt influence the behaviour and outcomes of a
RAG system when the user input, knowledge base, retrieval mechanism, and
underlying language model are held fixed?} \\
Answering it requires a formal model of the RAG system that identifies the
system prompt as a controlled parameter, and an experimental design that holds
every other source of variation fixed. \\
This thesis develops both for the simplest
retrieval-augmented architecture. In that architecture the system prompt reaches
the answer along one path, through the generator. A measured effect can
therefore be read as an effect on generation, and not as a change in what was
retrieved.


\section{Thesis Structure}

This thesis has seven chapters. Chapter 2 outlines why retrieval-augmented generation was developed and introduces the components of a
RAG system. Chapter 3 reviews the literature and identifies the research gap. Chapter 4 builds the formal model used here. It takes the probabilistic formulation of \citet{lewis2020rag} and extends it so that the system prompt becomes an explicit variable of the generator. It also instantiates the knowledge base with an established dataset and defines what counts as a good answer.
Chapter 5 describes the experimental setup. Chapter 6 presents and discusses the
results and states the limits of the design. Chapter 7 draws conclusions and outlines directions for future work.


\section*{Notation}
\label{sec:notation}
\addcontentsline{toc}{section}{Notation}

The notation follows \citet{lewis2020rag} for the RAG system and the intervention calculus of \citet{pearl2009causality} for the causal statements. 

\subsection*{The system and its components}

\begin{itemize}
    \item $q \in \mathcal{Q}$ denotes a user query, and $P_{\mathcal{Q}}$ the distribution over the query space $\mathcal{Q}$ from which queries are drawn.
    \item $D = \{d_1, d_2, \ldots, d_N\}$ denotes the knowledge base, where each $d_n$ is a document and $N$ is the number of documents.
    \item $S$ denotes a system prompt.
    \item $C$ denotes the context supplied to the language model: the passages it reads, in the order it reads them.
    \item $y = (y_1,\ldots,y_L)$ denotes the generated system output, a sequence of $L$ tokens; $y_{1:i-1}$ denotes the tokens produced before position $i$.
    \item $R$ and $G$ denote the retrieval and generation functions of the pipeline, and $\chi$ the chunking strategy by which documents are divided into retrievable units.
    \item $v \in \mathcal{V}$ denotes a token of the vocabulary $\mathcal{V}$, $u_v$ the score the generator assigns it, and $T > 0$ the decoding temperature.
\end{itemize}

\subsection*{The probabilistic model}

\begin{itemize}
    \item $z \in D$ denotes a single retrieved document, and $\mathcal{Z}_K(q,D)$ the $K$ documents with the highest retrieval score.
    \item $p_\eta(z \mid q,D)$ denotes the retrieval distribution with parameters $\eta$, and $p_\theta(y_i \mid q,C,y_{1:i-1},S)$ the generator with parameters $\theta$.
    \item $p_{\mathrm{RAG}}(y \mid q,D,S)$ denotes the distribution over outputs.
    \item $Y$ denotes the output viewed as a random variable, $M(Y)$ a measured property of it, and $\varepsilon$ the randomness introduced by decoding.
\end{itemize}

\subsection*{The system prompt and its variation}

\begin{itemize}
    \item $\Sigma = \{\sigma_1,\ldots,\sigma_m\}$ denotes the set of $m$ requirements a system prompt states.
    \item $\varphi$ denotes the wording in which those requirements are expressed, and $S = \rho(\Sigma,\varphi)$ denotes the resulting system prompt.
    \item $S_{b,0},\ldots,S_{b,3}$ denote the four prompt conditions of baseline $b \in \{A,B,C\}$, with $S_{b,0}$ the canonical form of that baseline, against which its three rewordings are compared. \textit{Baseline} always refers to one of the three published prompts $A$, $B$ and $C$.
\end{itemize}

\subsection*{Causal intervention}

\begin{itemize}
    \item $\mathrm{do}(S = S_{b,j})$ denotes the intervention that sets the system prompt to condition $S_{b,j}$, and $\mathrm{do}(C = c)$ the intervention that sets the context to $c$.
    \item $\tau_{b,j}$ denotes the average causal effect of condition $S_{b,j}$ relative to the canonical form $S_{b,0}$, taken over the query distribution $P_{\mathcal{Q}}$, and $\hat{\tau}_{b,j}$ its sample estimate. Both are defined separately for each measure $M$ and each evidence level $e$.
     \item $\mathbb{E}[\cdot]$ denotes expectation, and $\hat{\cdot}$ marks a quantity estimated from the sample.
\end{itemize}

\subsection*{Experimental design}

\begin{itemize}
    \item $c_{ie}$ denotes the context given to query $q_i$ at evidence level $e$, with $c_{i,\text{gold}}$ and $c_{i,\text{ret}}$ the two levels.
     \item $\mathcal{Q}_{\text{single}}$ and $\mathcal{Q}_{\text{multi}}$ denote the two query categories, and $N_Q$ the size of the query set.
    \item $A_q$ denotes the reference answers annotated for query $q$, against which correctness is measured.
    \item $I$ denotes the set of queries in one cell of the design, and $\bar{M}_{ij}$ the mean of a measure over the $n$ repetitions of query $i$ under prompt condition $j$.
    \item The indices are $i \in \{1,\ldots,N_Q\}$ over queries, $j \in \{0,\ldots,3\}$ over the wordings of a baseline, $e \in \{\text{gold},\text{ret}\}$ over evidence levels, and $r \in \{1,\ldots,n\}$ over the $n$ repetitions of each cell. The index $k \in \{1,\ldots,K\}$ is reserved for the rank of a retrieved passage.
\end{itemize}

\chapter{Preliminaries}



\section{Path towards RAG}
Large language models generate text from parametric knowledge alone, that is, from information encoded in their weights during training
\citep{wang2024knowledgemechanisms}. This leads to two problems. First, a model's
knowledge is fixed at the time its training data was collected. Information that appears later is unknown to the model unless the model is retrained or fine-tuned. Second, purely parametric models are prone to
hallucination: nothing in their design forces an output to trace back to a
verifiable source, so they can generate plausible-sounding statements that are unsupported \citep{ji2023hallucination}.

Both problems motivate giving a language model access to a knowledge base at inference time instead of relying only on what
is stored in its parameters \citep{lin2025llminference}. \citet{lewis2020rag}
formalised this idea as RAG. RAG combines a retriever, which selects relevant passages from a document
collection, with a pretrained generator, which conditions its output on those retrieved passages in addition to the input query.

\section{Components of RAG}

A RAG system has four components: the knowledge base the evidence comes from,
the retriever that selects relevant evidence for a given query, the generator
that produces a response from the query and the retrieved evidence, and the
system prompt that specifies how the generator should behave. Chapter 4 treats them formally.
Figure~\ref{fig:rag} shows how they fit together and which information passes between
them.

\begin{figure}[htbp]
\centering
\resizebox{\textwidth}{!}{% RAG system architecture -- style matched to Figure 4.1.
% Requires: \usetikzlibrary{arrows.meta,positioning,shapes.geometric,fit,backgrounds,calc}
\begin{tikzpicture}[
    box/.style   = {draw, rounded corners=2pt, minimum height=11mm, minimum width=30mm,
                    align=center, font=\small},
    store/.style = {draw, cylinder, shape border rotate=90, aspect=0.30,
                    minimum height=18mm, minimum width=28mm, align=center, font=\small},
    mod/.style   = {draw, dashed, rounded corners=2pt, inner sep=10pt},
    modlab/.style= {font=\small},
    stage/.style = {font=\small\scshape},
    ar/.style    = {-{Latex[length=2mm]}},
    arbi/.style  = {{Latex[length=2mm]}-{Latex[length=2mm]}},
    lb/.style    = {font=\scriptsize, inner sep=3pt, fill=white, align=center}
]

% ---- nodes -------------------------------------------------------------
\node[box]   (query)  at (0,0)      {User Query};

\node[box]   (search) at (4.6,0)    {Semantic Search /\\ Similarity Ranking};
\node[store] (vdb)    at (4.6,-3.2) {Vector Database\\ (Knowledge Base)};

\node[box]   (prompt) at (12.2,0)   {Prompt\\ Construction};
\node[box]   (llm)    at (16.2,0)   {Language Model};

\node[box]   (out)    at (20.6,0)   {Final Output\\ (Generated Answer)};

\node[box]   (sysp)   at (12.2,-3.2) {System Prompt};

% ---- module containers -------------------------------------------------
\begin{scope}[on background layer]
  \node[mod, fit=(search)(vdb)] (retbox) {};
  \node[mod, fit=(prompt)(llm)] (genbox) {};
\end{scope}

\node[modlab, anchor=south west] at ([yshift=3pt]retbox.north west) {Retrieval Module};
\node[modlab, anchor=south west] at ([yshift=3pt]genbox.north west) {Generation Module (LLM)};

\node[stage] at (4.6,3.5)  {I.\ The Retrieval Stage};
\node[stage] at (14.2,3.5) {II.\ The Generation Stage};

% ---- data flow ---------------------------------------------------------
\draw[ar]   (query)  -- (search);
\draw[arbi] (search) -- node[lb, midway, xshift=10mm] {embed /\\ rank} (vdb);

\draw[ar] ([yshift=-3mm]search.east)
          -- node[lb, pos=0.44] {Retrieved Context\\ (Top-$K$ Chunks)}
          ([yshift=-3mm]prompt.west);

\draw[ar] (query.north) -- (0,2.8)
          -- node[lb, pos=0.45] {User Query} (9.7,2.8)
          -- (9.7,0.3) -- ([yshift=3mm]prompt.west);

\draw[ar] (sysp) -- node[lb, midway, xshift=13mm] {instructions,\\ constraints, format} (prompt);

\draw[ar] (prompt) -- (llm);
\draw[ar] (llm)    -- (out);

\end{tikzpicture}
}
\caption{Architecture of a RAG system showing the four components introduced above. The retrieval stage matches the user query against a vector database and returns the top-$K$ chunks. The generation stage combines those chunks with the user query and the system prompt into a single prompt and produces the answer.}
\label{fig:rag}
\end{figure}

\subsection{Knowledge Base}

The knowledge base $D = \{d_1, d_2, \ldots, d_N\}$ is the collection of documents
the evidence comes from. In most RAG systems, documents are first split
into smaller passages, called chunks. 

An embedding model maps each chunk to a high-dimensional vector, so that
chunks with similar meaning lie close together in the embedding space.
Similarity between a query vector and a chunk vector is usually measured by the
cosine of the angle between them.

The embeddings are stored in a vector database, which supports fast
nearest-neighbour search over large collections of chunks. This thesis uses FAISS \citep{johnson2021faiss}, a library that is embedded directly in an application and gives fine-grained control over the index structure, because it supports the exact search that Section~\ref{sec:procedure} requires.
\subsection{Retriever}

Given a user query $q$ and a knowledge base $D$, the retriever selects the chunks
that are relevant for answering $q$. Chapter 4 formalises this as a probability
distribution $p_\eta(z \mid q,D)$ over the documents of the knowledge base.
\subsubsection{Dense Retrieval}

Dense retrieval embeds the query with the same embedding model used for the
knowledge base and ranks chunks by the similarity between the query vector and
each chunk vector. Because the embedding is learned, two passages that express
the same content in different words lie close together in the embedding space. A query can therefore match
a chunk with which it shares no vocabulary.

\subsubsection{Retrieval Depth}

The top-$K$ ranked chunks form the retrieved context. The
experiment of Chapter 5 fixes the retriever, the chunking and $K$, and controls
the order of the passages.

\subsection{Generator}
\label{sec:generator}

The generator is the large language model that produces the final response. It generates text by autoregressive decoding. Its parameters, written
$\theta$ in Chapter 4, encode what the model learned during training. This
encoded knowledge is often called parametric memory.

The three inputs (the system prompt, the retrieved context and the query) reach the model in separate message roles. The system prompt is
supplied in the system role, and the retrieved context and the query in the user
role. The model's chat template then assembles the messages into a single token sequence, with the system prompt first, so that it conditions the processing of every subsequent token. Chapter 4 formalises the generator as
$p_\theta(y_i \mid q,C,y_{1:i-1},S)$, the probability of the next token given the
query, the retrieved context, the tokens already produced and the system prompt.
Changing any one of these inputs can change the output, which is why the
experiment varies only $S$.

\subsubsection{Decoding and the Temperature Parameter}

Generation is not deterministic by default. At each step the model produces a
score $u_v$ for every token $v$ in the vocabulary $\mathcal{V}$. The softmax function, with a temperature parameter $T > 0$, turns these scores into a
probability distribution:

\begin{equation}
\label{eq:temperature}
p(v) = \frac{\exp(u_v / T)}{\sum_{v' \in \mathcal{V}} \exp(u_{v'} / T)}.
\end{equation}

The temperature rescales the scores before they are normalised, and it controls
how much probability mass is left for tokens the model considers less likely
\citep{holtzman2020degeneration}. As $T$ approaches zero, the distribution
concentrates on the highest-scoring token, and decoding becomes effectively
deterministic. As $T$ grows, the distribution flattens.

This thesis holds the temperature $T$ of Equation~\eqref{eq:temperature} fixed at the API default of $1.0$ across all experimental conditions. The comparison between prompts requires only
that $T$ be constant. 
\subsection{System Prompt}

The system prompt $S$ specifies the instructions, constraints and output
requirements the generator should follow. It has two functions. The first is to define the role the RAG system occupies, the audience it addresses,
and the format its output takes. The second is to constrain the system's behaviour; in a RAG system, the central constraint is the grounding constraint. A grounding constraint restricts the model to the
retrieved context and specifies what to do when that context does not suffice.


\section{Advanced RAG Architectures}
\label{sec:advanced}

The architecture described so far retrieves once per query and generates once
from what was retrieved. Recent variants depart from this procedure and can be grouped by the stage of the pipeline they modify.

\begin{description}
    \item[Conversational RAG.] Storing prior turns of a dialogue as state, so that follow-up questions can be answered \citep{liu2024chatqa}.
    \item[Pre-retrieval optimisation.] Rewriting or expanding the query before it reaches the retriever. 
    \item[Routing and adaptive retrieval.] Deciding where and how much to retrieve. 
    \item[Self-correction and reflection.] Evaluating retrieved evidence before or during generation. Self-RAG instead trains the model to decide for itself when to retrieve and to critique its own draft output, so that retrieval and generation are interleaved \citep{asai2024selfrag}.
    \item[Agentic RAG.] Treating retrieval as one action among several available to an agent. A meta-agent delegates to document-specific or tool-specific agents and combines their results, which turns a fixed pipeline into a workflow whose steps are chosen at run time \citep{singh2025agentic}.
\end{description}

Each variant
changes the retrieval function $R$, the generation function $G$, or the number of times they are applied. All of them, however, are still built from a retrieval step and a generation step. 

This thesis implements only the simplest architecture and none of the variants above.
The reason is methodological: the objective is to estimate the causal effect of
the system prompt on the output, which Chapter 4 formalises as the intervention
$\mathrm{do}(S = S_{b,j})$. Because the prompt is set by intervention and
everything else is held fixed, its total effect on the output can be estimated
in any architecture. The simple architecture, however, makes the effect interpretable. In this architecture, the system prompt reaches the answer along a single path, through the
generator, so the effect is an effect on generation. In the variants above the
prompt, or the text it shapes, can also change what is retrieved, how often, or from where. 
\chapter{Literature Review}

This chapter reviews two
bodies of literature relevant to how the system prompt affects such a system (prompt sensitivity in language models, and retrieval-augmented generation itself) and then identifies the gap between them.

\section{Prompt Sensitivity and Instruction-Following in Language Models}
\label{sec:promptlit}

The first body of literature studies language models with no retrieval
component. \citet{cheng2026systemprompts} vary
the system prompt given to instruction-tuned models for code generation, crossing
prompt specificity and prompting strategy with model, programming language and
decoding settings. They report that the system prompt influences the quality of
generated code independently of the task description. \citet{mu2025robustness} examine how reliably models follow a system prompt when the
user turn contains instructions conflicting with it, and report that the system
prompt does not always prevail. \citet{zheng2024helpful} assign personas through the system prompt across a large set of roles and find that doing so does not reliably improve accuracy over a system prompt with no persona. Together, these three studies establish the system prompt as a channel with measurable influence over behaviour, and show that the influence is not guaranteed.

A larger body of research varies the user-side prompt.
\citet{zhuo2024prosa} propose ProSA, a framework that measures sensitivity to instance-level variation in the task instruction and relates it to decoding confidence. They report that larger models are more robust to such variation and that few-shot examples reduce sensitivity.
\citet{salinas2024butterfly} perturb task prompts across text classification
tasks (adding a trailing space, requesting the output in XML, or prefixing a
jailbreak) and find that even such small perturbations change the labels the model
assigns. \citet{sclar2024quantifying} vary only the formatting of few-shot task
templates (separators, casing and spacing, none of which carries propositional
content) and report differences of up to 76 accuracy points between formats a
reader would judge equivalent. \citet{hua2025flaw} evaluate twelve task templates
across several benchmarks and argue that much of the measured sensitivity is
produced by the evaluation method (log-likelihood scoring and rigid answer
matching).

\citet{ying2024intuitive} study conflict
between a model's parametric memory and the evidence supplied in its context, and
classify models as intuitive or dependent according to which source they favour. 

Two conclusions follow. The first, supported directly by the three system-prompt studies \citep{cheng2026systemprompts,mu2025robustness,zheng2024helpful}, is that the system prompt is a behavioural control variable whose effect can be isolated and measured in a properly controlled setup.
The second is that model behaviour is sensitive to changes in the prompt that preserve its meaning. This holds whether the change is a rephrasing of the instruction \citep{zhuo2024prosa}, a perturbation as small as a trailing space \citep{salinas2024butterfly}, or a change of formatting alone \citep{sclar2024quantifying}. The apparent size of such an effect, however, depends on how it is measured \citep{hua2025flaw}. This second conclusion is established
for the user-side prompt, and its transfer to the system prompt is an assumption rather than an empirical finding. The assumption is reasonable, since both texts reach the
same decoder as tokens in the same context window, but it is untested.

\section{Retrieval-Augmented Generation: Architecture, Evaluation, and Failure Modes}

The second body of literature concerns retrieval-augmented generation itself.
\citet{lewis2020rag} introduce RAG, combining a
retriever with a pretrained sequence-to-sequence generator, and show that a generator conditioned on retrieved evidence outperforms purely
parametric models on knowledge-intensive tasks. Their probabilistic formulation,
in which retrieval is a distribution over documents and the retrieved document is
marginalised out, is the direct basis for the model formalised in Chapter 4.
\citet{gao2023survey} survey the development of RAG systems through Naive,
Advanced and Modular paradigms. Their taxonomy treats retrieval and generation as
separable modules, which allows this thesis to hold one fixed while
studying an input to the other. Two later developments
extend this taxonomy in directions this thesis deliberately does not follow.
\citet{edge2024graphrag} restructure the retrieval side into a knowledge graph,
grouped into thematic segments and pre-summarised by an LLM, so that a query is
answered by traversing entities and summaries. \citet{singh2025agentic} and
\citet{jin2025searchr1} restructure the generation side into an agent loop, in
which the model itself decides when and what to retrieve, calling $R$ repeatedly
and adaptively. Both directions relax exactly the components Chapter 4
holds fixed: $R$ in the first case, and the single-call structure of $G$ in the
second. This thesis therefore retains \citeauthor{gao2023survey}'s simplest,
Naive-RAG form, for the reason given in Section~\ref{sec:advanced}.

\citet{es2023ragas} propose Ragas, a reference-free evaluation framework for RAG
pipelines built around three metrics that need no ground-truth human annotation:
faithfulness, answer relevance and context relevance. This thesis adopts faithfulness and answer relevance. The third, context
relevance, is a retrieval-side measure and is excluded for the reasons given in
Section~\ref{sec:metrics}. \citet{niu2024ragtruth} introduce RAGTruth, a manually annotated corpus of
word-level hallucinations produced by RAG systems, and benchmark
hallucination-detection methods across multiple language models. Their corpus
shows that RAG outputs contain claims the retrieved context does not support,
which is the failure the faithfulness measure of this thesis captures.

Prompts have been compared within RAG research, but as a design choice rather
than as the object of study \citet{mao2024fitrag}.

\section{Synthesis and Research Gap}

The first body of work shows two things: the system prompt is a controllable way of shaping language model behaviour, and model behaviour is sensitive to rewording that leaves the meaning of a prompt unchanged. Both are established without retrieval. The second body of work establishes the architecture, evaluation vocabulary and known failure modes of RAG systems. Where it varies the prompt, it varies what the prompt asks for, in order to choose a better one.

The gap lies between the two. No study reviewed here holds a RAG pipeline fixed and varies only the wording of its system prompt while keeping its requirements constant. This thesis addresses this gap.

\chapter{Formal Model of a RAG System}

Chapter 3 identified the research gap: prompt sensitivity has so far been studied only outside retrieval, and RAG research varies prompts only to choose between them, not to measure the effect of their wording. Closing that gap requires a model of the system precise enough to isolate the system prompt. This chapter builds that model on the probabilistic formulation of RAG given by \citet{lewis2020rag}, adapts it to a generator that reads all retrieved passages at once, and makes the system prompt an explicit variable of the generation process.

\section{Modelling Objective}



The model can be read as a chain,
\[
(q,D) \;\longrightarrow\; C \;\longrightarrow\; p_\theta(y_i \mid q,C,y_{1:i-1},S) \;\longrightarrow\; p_{\mathrm{RAG}}(y \mid q,D,S),
\]
in which retrieval produces the context $C$ and the system prompt $S$ enters at the generation stage.

\section{Probabilistic Retrieval}

Retrieval is represented as a probability distribution over documents:

\begin{equation}
\label{eq:retrieval}
p_\eta(z \mid q,D).
\end{equation}

Equation~\eqref{eq:retrieval} denotes the probability that the retriever assigns to document $z$, given the query $q$ and the knowledge base $D$. \citet{lewis2020rag} obtain it by normalising the retrieval scores with a softmax over the documents.  The formulation keeps only the $K$ documents with the highest retrieval probability,

\begin{equation}
\label{eq:topk}
\mathcal{Z}_K(q,D) = \operatorname*{TopK}_{z \in D} \; p_\eta(z \mid q,D),
\end{equation}

and the retrieved context passed to the generator is the resulting sequence of passages,

\begin{equation}
\label{eq:context}
C = (z_1,\ldots,z_K), \qquad z_k \in \mathcal{Z}_K(q,D).
\end{equation}

$C$ is an ordered sequence, not a set, because the generator reads the passages in a particular order and that order can affect the answer \citep{liu2024lostmiddle}. 

\section{Prompt-Conditioned Generation}

\citet{lewis2020rag} write the generator at the level of a single token as $p_\theta(y_i \mid q,z,y_{1:i-1})$. It is the probability of the next token given the query, one retrieved document, and the tokens already produced. With the system prompt added as a further conditioning variable, this becomes $p_\theta(y_i \mid q,z,y_{1:i-1},S)$. Because generation is autoregressive, the probability of a complete output under document $z$ follows by multiplication:

\begin{equation}
\label{eq:sequence}
p_\theta(y \mid q,z,S) = \prod_{i=1}^{L} p_\theta(y_i \mid q,z,y_{1:i-1},S).
\end{equation}

\section{From One Document to the Whole Context}
\label{sec:ragdist}

Equation~\eqref{eq:sequence} gives the probability of an output under a single document $z$. Since the system does not know which retrieved document supports the answer, \citet{lewis2020rag} average over the retrieved documents, weighting each by its retrieval probability. They give two ways of doing this.

\subsection{RAG-Sequence}

RAG-Sequence assumes that a single retrieved document is responsible for the whole answer. The probability of an output is the weighted average, over documents, of the probability that document would produce that entire sequence:

\begin{equation}
\label{eq:ragseq}
p_{\text{RAG-Seq}}(y \mid q,D,S) \approx \sum_{z \in \mathcal{Z}_K(q,D)} p_\eta(z \mid q,D) \prod_{i=1}^{L} p_\theta(y_i \mid q,z,y_{1:i-1},S).
\end{equation}

\subsection{RAG-Token}

RAG-Token makes the weaker assumption that different parts of an answer may come from different documents. The average over documents is taken separately at every token:

\begin{equation}
\label{eq:ragtok}
p_{\text{RAG-Token}}(y \mid q,D,S) \;\approx\; \prod_{i=1}^{L} \sum_{z \in \mathcal{Z}_K(q,D)} p_\eta(z \mid q,D) \, p_\theta(y_i \mid q,z,y_{1:i-1},S).
\end{equation}

Both forms run the generator once per document. The system studied here does not. Like the VANILLA setting of \citet{gao2023alce}, it concatenates the $K$ passages into one prompt and runs the generator once, so that the model reads all of them together and decides for itself which to use. The average over documents is then no longer computed explicitly. It is replaced by a generator that conditions on the whole context $C$ of Equation~\eqref{eq:context}:

\begin{equation}
\label{eq:token}
p_\theta(y_i \mid q,C,y_{1:i-1},S).
\end{equation}

Equation~\eqref{eq:token} defines the generator used in this thesis. It is here that the thesis departs from the original formulation in two ways: the model reads the whole context rather than one document, and the system prompt becomes part of what the generator conditions on. The distribution over complete outputs is

\begin{equation}
\label{eq:ragdist}
p_{\mathrm{RAG}}(y \mid q,D,S) \;=\; \prod_{i=1}^{L} p_\theta\bigl(y_i \mid q,C(q,D),y_{1:i-1},S\bigr),
\end{equation}

where $C(q,D)$ is the context retrieved for $q$ from $D$. Because the retriever is deterministic, the context is a function of the query and the knowledge base. The only randomness left in Equation~\eqref{eq:ragdist} is that of decoding. 

\section{The System Prompt: Requirements and Wording}
\label{sec:promptreq}

The system prompt enters Equation~\eqref{eq:token} as a single symbol. It is, however, a text with internal structure, and the experiment varies its wording deliberately. A prompt states a set of requirements in a particular wording. Separating the two gives

\begin{equation}
\label{eq:prompt}
S = \rho(\Sigma,\varphi), \qquad \Sigma = \{\sigma_1,\ldots,\sigma_m\}.
\end{equation}

This formulation separates what a prompt asks for from how it is worded. 

Three requirement sets are used. Each was taken from a prompt published in retrieval-augmented generation research and is reproduced in Section~\ref{sec:conditions}. $\Sigma_A$ is the response-generation prompt of the RAGBench benchmark \citep{friel2024ragbench}, which states a single requirement. $\Sigma_B$ is the VANILLA instruction of \citet{gao2023alce}, which comprises about ten clauses in four families (accuracy and concision, source restriction, citation, and tone). Of these, only citation can be checked by rule. $\Sigma_C$ is the instruction of \citet{mao2024fitrag}, which states three.

The second component is the wording. For each requirement set, four wordings are written that are intended to hold $\Sigma$ fixed and change only $\varphi$. The lexical and syntactic dimensions are taken from three prior studies. \citet{wahle2024paraphrase} vary prompts by linguistically defined paraphrase type, and group those types into families, among them lexicon and syntax. \citet{zhan2024lexical} examine the lexical dimension by replacing individual words with semantically similar ones, and report that models are over-sensitive to variations undetectable to a reader. \citet{kostic2026samemeaning} separate a lexical pipeline built on synonym substitution from a syntactic one built on dependency parsing. They find that the two behave differently: lexical perturbation degrades performance consistently, while syntactic perturbation has more heterogeneous effects and occasionally improves results. The third dimension, layout, follows \citet{sclar2024quantifying}, who change only the formatting of a prompt.

This thesis follows their approach of varying the system prompt along linguistically defined, meaning-preserving dimensions one at a time. Writing $S_{b,j}$ for the $j$-th wording of baseline $b$:
\begin{description}
    \item[$S_{b,0} = \rho(\Sigma_b,\varphi_0)$, canonical.] The published form, unchanged.
    \item[$S_{b,1} = \rho(\Sigma_b,\varphi_1)$, lexical.] Synonym substitution with clause order and voice held fixed.
    \item[$S_{b,2} = \rho(\Sigma_b,\varphi_2)$, syntactic.] Clause reordering and conversion from active to passive voice.
    \item[$S_{b,3} = \rho(\Sigma_b,\varphi_3)$, format.] The same wording laid out as a list rather than as prose.

\end{description}

Baseline $A$ states the fewest requirements and is short enough to be given in full as an example. The four conditions read:

\begin{lstlisting}[style=promptbox,caption={The four conditions of baseline $A$ from \citet{friel2024ragbench}.},label={lst:baselineA}]
A0  Use the following pieces of context to answer the question.

A1  Employ the following excerpts of material to address the query.

A2  The question is to be answered using the context pieces below.

A3  CONTEXT: the following pieces
    TASK: use them to answer the question
\end{lstlisting}

The variants for baselines $B$ and $C$ are written on the same principle and are longer in proportion to their requirement sets. Section~\ref{sec:conditions} gives them in full, together with the places where a rewording does not fully meet this principle.

\section{The System Prompt as an Intervention}
\label{sec:intervention}

\begin{figure}[htbp]
\centering
\resizebox{\textwidth}{!}{%
\begin{tikzpicture}[
    box/.style   = {draw, rounded corners=2pt, minimum height=9mm, minimum width=21mm,
                    align=center, font=\small},
    tx/.style    = {draw, very thick, rounded corners=2pt, minimum height=9mm,
                    minimum width=21mm, align=center, font=\small},
    ar/.style    = {-{Latex[length=2mm]}},
    lb/.style    = {font=\scriptsize, inner sep=2pt, fill=white, align=center}
]
\node[box] (q)    at (0,1.4)    {$q$};
\node[box] (d)    at (0,-1.4)   {$D$};
\node[box] (retr) at (3.9,-1.4) {Retriever $\eta$};
\node[box] (z)    at (7.6,-1.4) {$C$};
\node[tx]  (s)    at (7.6,2.4)  {$S$};
\node[box] (gen)  at (11.3,0)   {Generator $\theta$};
\node[tx]  (y)    at (14.9,0)   {$y$};
\node[box] (m)    at (18.4,0)   {$M(Y)$};

\draw[ar] (q) -- (retr);
\draw[ar] (d) -- (retr);
\draw[ar] (retr) -- (z);
\draw[ar] (z) -- (gen);
\draw[ar] (q) -- (gen);
\draw[ar, very thick] (s) -- (gen);
\draw[ar] (gen) -- (y);
\draw[ar] (y) -- (m);
\end{tikzpicture}%
}
\caption{The causal structure of the system. The system prompt $S$ and the output $y$ are drawn in bold. $S$ enters the generator and has no arrow into the retriever or the context $C$, so its effect on $y$ runs through generation alone.}
\label{fig:causal}
\end{figure}

For the effect of the prompt to be interpretable, the following are held constant across conditions. Each query is used unchanged in every condition in which it is run, so that every comparison is made within a query rather than between queries, and the knowledge base is fixed at $D = D_0$. The retriever and its parameters are fixed at $\eta = \eta_0$, and the language model at $\theta = \theta_0$.

Decoding is also stochastic, so the same configuration can produce different outputs from one run to the next. The term $\varepsilon$ captures that variation:

\begin{equation}
\label{eq:noise}
Y = G(q,C,S;\theta,\varepsilon),
\end{equation}

which is why each condition is repeated several times.

\section{The Causal Effect}

Setting the system prompt to a fixed value is an intervention, written $\mathrm{do}(S = S_{b,j})$ following \citet{pearl2009causality}. The notation separates \textit{setting} a variable from merely \textit{observing} it. For a query $q$ under prompt condition $S_{b,j}$, the output of the system is the random variable
\begin{equation}
\label{eq:potential}
Y \sim p_{\mathrm{RAG}}(y \mid q,D,S_{b,j}).
\end{equation}

For one query the effect of the prompt is the difference in the expected value of the measured property $M(Y)$ of Equation~\eqref{eq:potential} between the two conditions. The research question concerns the behaviour of the system on the queries it will actually encounter. The effect therefore has to be defined as an average over the distribution from which queries are drawn. Writing $P_{\mathcal{Q}}$ for that distribution over the query space $\mathcal{Q}$, the causal effect of prompt condition $S_{b,j}$ relative to the canonical form $S_{b,0}$ of the same baseline is
\begin{equation}
\label{eq:estimand}
\tau_{b,j} \;=\; \mathbb{E}_{q \sim P_{\mathcal{Q}}}\Bigl[\; \mathbb{E}_{Y \sim p_{\mathrm{RAG}}(\cdot \mid q,D,S_{b,j})}\bigl[M(Y)\bigr] \;-\; \mathbb{E}_{Y \sim p_{\mathrm{RAG}}(\cdot \mid q,D,S_{b,0})}\bigl[M(Y)\bigr] \;\Bigr],
\end{equation}

which connects the probability model of Section~\ref{sec:ragdist} directly to the quantity Chapter 5 estimates. 

Equation~\eqref{eq:estimand} is an average treatment effect in the standard sense: the query plays the role of the unit, the prompt condition the role of the treatment, and $\tau_{b,j}$ the mean over units of the difference between the two treatments. The expectations are nested: the inner expectation averages over the randomness of decoding at a fixed query and a fixed context, while the outer averages over the queries themselves.

The distribution $P_{\mathcal{Q}}$ is not known and cannot be written down. Chapter 5 replaces it with the empirical distribution of a query set. The choice of data therefore determines how far the results generalise, which is why Section~\ref{sec:knowledgebase} draws the queries from an established question-answer dataset.

The choice of $M$ is left open here. Section~\ref{sec:goodanswer} describes what that property has to capture, and Section~\ref{sec:metrics} fixes it.

\section{Grounding the Knowledge Base in an Established Dataset}
\label{sec:knowledgebase}

$D$ and the query set are instantiated from an established question-answer dataset, MS MARCO \citep{nguyen2016msmarco}. The corpus is public, so a third party can repeat the experiment, and its answers are scored with established metrics.

The collection is built from queries sampled from Bing search logs, with candidate passages retrieved for each query and free-form answers written by human judges. MS MARCO provides the queries $q$, the documents that populate $D$, and, through its annotations, the reference answers $A_q$ and the passages the annotators marked as supporting them.

What it does not provide is a definition of a good answer. 

\section{What Counts as a Good Answer}
\label{sec:goodanswer}
Answer quality in a retrieval-augmented system comprises several properties that can fail independently of one another. The first is groundedness: an answer should rest on the documents the retriever actually supplied. The second is relevance: an answer should address the question that was asked, and address it directly. The third is correctness: an answer should agree with the reference answer recorded in the dataset, which is a different demand from the first and is treated separately here for the reason given below. The fourth is instruction adherence: whatever the system prompt asks for should be visible in the output itself, whether in length, in form, or in the handling of missing evidence. Instruction adherence is the property through which a prompt is meant to act, although a prompt can also change an answer in ways it did not ask for.

Groundedness and correctness are kept apart deliberately, because collapsing them would conceal the failure this thesis is concerned with. The retrieved context is neither complete nor guaranteed to be free of error. An answer can therefore be correct while ignoring the evidence entirely, drawing instead on what the generator already holds in its parameters. It can also use the evidence faithfully and still be wrong, because the evidence was wrong. A high correctness score alone would reveal nothing about whether retrieval was used at all, and the use of retrieval is exactly what a system prompt is supposed to govern.

\chapter{Experimental Setup}
 
\section{Research Question and Hypothesis}
\label{sec:goal}

The experiment operationalises the research question of Chapter 1 as follows: \textit{does a RAG system behave differently when the wording of its system prompt is changed while the meaning of the prompt stays the same?}

Wording refers to how the prompt is written: which words are used, how the sentences are built, and how the text is laid out. Meaning refers to the set of requirements $\Sigma$ that the prompt states, in the sense of Equation~\eqref{eq:prompt}. Every condition of the experiment keeps $\Sigma$ fixed and changes only the wording $\varphi$. If the system responded only to what a prompt requires and not to how it is written, all wordings of one prompt would produce the same distribution of answers. Whether they do is the empirical question. Stated as hypotheses, it reads:

\begin{description}
    \item[$H_0$.] System prompts that state the same requirements in different wordings produce the same distribution of answer quality: $\tau_{b,j} = 0$ for every rewording $j$, in every baseline, on every measure and at both evidence levels.
    \item[$H_1$.] At least one rewording changes answer quality: $\tau_{b,j} \neq 0$ for some rewording in some baseline on some measure at some evidence level.
\end{description}


\section{Overview of the Experiment}
\label{sec:overview}

The system under test is the simple RAG architecture of Chapter 2. Table~\ref{tab:fixed} lists the components of that system and the value each one is fixed at.

\begin{table}[htbp]
\centering
\caption{The components of the system under test and their fixed values.}
\label{tab:fixed}
\begin{tabular}{l>{\raggedright\arraybackslash}p{0.6\textwidth}}
\toprule
Component & Fixed value \\
\midrule
Knowledge base $D$ & passages of the 12,467 MS MARCO v2.1 validation queries with a well-formed answer, $N = 122{,}678$ \\
Chunking $\chi$ & one passage per chunk, as published \\
Embedding model & \texttt{all-mpnet-base-v2} \citep{reimers2019sbert}, mean pooling, L2-normalised \\
Vector index & FAISS flat inner-product index, exact search \\
Retrieval depth $K$ & 10 \\
Generator $\theta$ & \texttt{gpt-4o-mini} \citep{openai2024gpt4omini} via the chat completions API \\
Decoding & temperature $1.0$ (the API default) \\
Judge & \texttt{gpt-4.1-mini} \citep{openai2025gpt41}, temperature $0$, fixed prompts \\
Query set & 100 MS MARCO queries, seed 20260904 \\
Repetitions $n$ & 3 per cell \\
\bottomrule
\end{tabular}
\end{table}

The experiment has three stages, of which the first two collect data. In the first stage, retrieval is run once for every query and its result is written to disk. In the second stage, every prompt condition is applied to every query using the stored retrieval results, and the answers are written to disk. In the third stage, the answers are scored and the prompt effect is estimated.

\section{System-Prompt Conditions}
\label{sec:conditions}
Baseline $A$ is the response-generation prompt of \citet{friel2024ragbench}, reproduced in Listing~\ref{lst:minimal}. It states one requirement: answer from the supplied context.

\begin{lstlisting}[style=promptbox,caption={Baseline $A$: the response-generation prompt of \citet{friel2024ragbench}.},label={lst:minimal}]
Use the following pieces of context to answer the question.
{documents}
Question: {question}
\end{lstlisting}

Baseline $B$ is the VANILLA instruction of \citet{gao2023alce}. It constrains accuracy, engagement, concision and register, it fixes the syntax and the density of citations, and it warns that some of the supplied results may be irrelevant.

\begin{lstlisting}[style=promptbox,caption={Baseline $B$: the VANILLA instruction of \citet{gao2023alce}, reproduced from their Table 23.},label={lst:alce}]
Write an accurate, engaging, and concise answer for the given
question using only the provided search results (some of which
might be irrelevant) and cite them properly. Use an unbiased and
journalistic tone. Always cite for any factual claim. When citing
several search results, use [1][2][3]. Cite at least one document
and at most three documents in each sentence. If multiple documents
support the sentence, only cite a minimum sufficient subset of the
documents.
\end{lstlisting}

Baseline $C$ is the instruction of \citet{mao2024fitrag}. Its authors describe it as consisting of three parts: the first directs the model to the retrieved passages, the second asks it to read the question carefully, and the third asks for an answer together with a justification. 

\begin{lstlisting}[style=promptbox,caption={Baseline $C$: the comprehensive instruction of \citet{mao2024fitrag}.},label={lst:fitrag}]
Refer to the passage below and answer the following question.
Make sure you fully understand the meaning of the question and
passages. Then give the answer and explain why you choose this
answer.
Passages: {passages}
Question: {question}
\end{lstlisting}

The published prompts contain placeholders for the documents and the question. In this experiment the instruction text alone forms the system prompt $S$ and is sent in the system role of the chat interface. The passages and the question are sent in the user role, in the fixed template of Listing~\ref{lst:userturn}. The template is the same in every condition, so the only text that differs between conditions is $S$. Passages are numbered from one, so that citations can refer to them and the position of a passage can be recovered from the answer.

\begin{lstlisting}[style=promptbox,caption={The user turn. It is identical in every condition.},label={lst:userturn}]
Passages:
[1] <first passage>
[2] <second passage>
...

Question: <the query>
Answer:
\end{lstlisting}



A rewording that changed how strongly an instruction was stated, for example by turning \textit{always cite} into \textit{please try to cite}, would change what the prompt requires. Such a change is a change of $\Sigma$, and a difference in behaviour caused by it could not be read as an effect of wording. The rewordings were written to avoid such changes. In four places, however, they come close. These are listed here so that the results can be interpreted with them in mind:
\begin{itemize}
    \item C1 asks the model to \textit{justify why you select} its answer where C0 asks it to \textit{explain why you choose} it. \textit{Justify} may ask for slightly more than \textit{explain}.
    \item C3 adds \textit{and nothing else} to the source line, which states an exclusivity that C0 leaves implicit.
    \item B2 and C2 turn several instructions into plain statements (\textit{A citation is always given}, \textit{The answer is then given}) rather than obligations (\textit{is to be given}), which may state them less strongly.
    \item B3 and C3 keep some items as imperatives (\textit{cite the search results properly}, \textit{give the answer}), so the list format does not eliminate every imperative.
\end{itemize}

The four wordings of baseline $A$ are given in Listing~\ref{lst:baselineA} in Chapter 4. Listings~\ref{lst:baselineB} and \ref{lst:baselineC} give the three rewordings of baselines $B$ and $C$; their canonical forms are Listings~\ref{lst:alce} and \ref{lst:fitrag} above.

\begin{lstlisting}[style=promptbox,caption={The lexical ($B1$), syntactic ($B2$) and format ($B3$) wordings of baseline $B$.},label={lst:baselineB}]
B1  Compose a precise, compelling, and succinct reply for the posed
    query drawing only on the supplied search results (a number of
    which may be immaterial) and reference them correctly. Adopt an
    impartial and journalistic register. Invariably reference every
    factual assertion. When referencing several search results,
    employ [1][2][3]. Reference at minimum one document and at
    maximum three documents per sentence. If several documents
    corroborate the sentence, only reference a minimal adequate
    subset.

B2  For the given question an accurate, engaging, and concise answer
    is to be written from the provided search results alone, some
    of which might be irrelevant, and cited properly. An unbiased,
    journalistic tone is to be used. A citation is always given for
    any factual claim. Where several results are cited, [1][2][3]
    is used. In each sentence at least one and at most three
    documents are cited. Where multiple documents support a
    sentence, only a minimum sufficient subset is cited.

B3  ANSWER
    - accurate, engaging, and concise, for the given question
    - source: only the provided search results
    - note: some of the provided results might be irrelevant

    CITATIONS
    - cite the search results properly
    - always cite for any factual claim
    - when citing several search results: [1][2][3]
    - in each sentence: at least one document, at most three documents
    - if multiple documents support the sentence: only a minimum
      sufficient subset

    TONE
    - unbiased and journalistic
\end{lstlisting}

\begin{lstlisting}[style=promptbox,caption={The lexical ($C1$), syntactic ($C2$) and format ($C3$) wordings of baseline $C$.},label={lst:baselineC}]
C1  Consult the excerpt below and respond to the following query.
    Ensure you completely grasp the sense of the query and excerpts.
    Then supply the response and justify why you select this
    response.

C2  The following question is to be answered from the passage below.
    The meaning of the question and passages is first understood.
    The answer is then given, with an explanation of why it is
    chosen.

C3  SOURCE: the passage below, and nothing else
    READ: fully understand the meaning of the question and passages
    ANSWER: give the answer to the following question
    EXPLAIN: say why you choose this answer
\end{lstlisting}

Prompt length is held nearly constant within a baseline, so that a difference between two wordings cannot be an effect of how much text the prompt contains. Table~\ref{tab:lengths} gives the word counts. No wording departs from its canonical form by more than one tenth.
\begin{table}[htbp]
\centering
\caption{Word count of every condition, with the deviation from the canonical form of its baseline.}
\label{tab:lengths}
\begin{tabular}{lrrrr}
\toprule
Baseline & canonical & lexical & syntactic & format \\
\midrule
$A$ & 10 & 10 ($\pm 0\%$) & 11 ($+10\%$) & 11 ($+10\%$) \\
$B$ & 74 & 71 ($-4\%$) & 80 ($+8\%$) & 77 ($+4\%$) \\
$C$ & 33 & 32 ($-3\%$) & 34 ($+3\%$) & 32 ($-3\%$) \\
\bottomrule
\end{tabular}
\end{table}

This yields twelve conditions: three baselines in four wordings each. They are written $S_{b,j}$ with $b \in \{A,B,C\}$ and $j \in \{0,1,2,3\}$. Comparing $S_{B,1}$ with $S_{A,0}$ would mix a change of wording with a change of what is asked for. Hence, each rewording is compared only with the canonical form of its own baseline.

\section{Variables and Controlled Factors}
\label{sec:variables}

\paragraph{Independent variable.} The wording $j$ of the system prompt for each baseline $b$. 

\paragraph{Dependent variable.} The quality of the answer, as scored by the measures of Section~\ref{sec:metrics}: faithfulness, answer relevance, correctness and instruction adherence, with correctness measured by three metrics: ROUGE-L, BLEU-1 and BERTScore. Each of the six is used in turn as the measured property $M$ of Equation~\eqref{eq:estimand}, so the effect of a rewording is estimated six times. All six take part in the test of $H_0$. They are reported separately.

\paragraph{Controlled factor.}  The context $C$ passed to the generator is set directly, as the intervention $\mathrm{do}(C = c_{ie})$ of Section~\ref{sec:ragdist}, at two evidence levels $e$:
\begin{description}
    \item[$c_{i,\text{gold}}$.] Only the passages the annotation marks as supporting the answer. The evidence needed for the answer is guaranteed to be present.
    \item[$c_{i,\text{ret}}$.] The $K$ passages the fixed retriever actually returns. 
\end{description}
The two levels differ in more than the quality of the evidence. A gold context holds one or a few passages, a retrieved context always ten.


\paragraph{Held constant.} Everything in Table~\ref{tab:fixed}. Queries and passages are embedded separately by the same model and compared by their vectors, as in the bi-encoder of \citet{karpukhin2020dpr}, but with one shared encoder for both sides \citep{reimers2019sbert}. Because the embedding model, the index, the chunking and the depth are all fixed, the set $\mathcal{Z}_K(q,D)$ of Equation~\eqref{eq:topk} depends on the query alone and never on the system prompt. 
\paragraph{Covariates.} Two properties of each query are recorded: the category (single-passage or multi-passage) and MS MARCO's intent type. Both are defined in Section~\ref{sec:corpus}. They are balanced in the sample so that it covers different kinds of query. The analysis of this thesis does not break the effects down by them.

\paragraph{The design.} An observation is one answer, written
\begin{equation}
\label{eq:design}
y_{ijer} \;\sim\; p_{\mathrm{RAG}}\bigl(y \;\big|\; q_i,\ \mathrm{do}(C = c_{ie}),\ \mathrm{do}(S = S_{b,j})\bigr), \qquad r = 1,\ldots,n,
\end{equation}
for query $i$, prompt condition $S_{b,j}$, evidence level $e$ and repetition $r$. The twelve prompt conditions and the two evidence levels are fully crossed, which gives 24 cells. With 100 queries and $n = 3$ repetitions this would be 7,200 generations. Queries that fail the retrieval gate of Section~\ref{sec:procedure} are not run at the retrieved level, so fewer answers are generated (Chapter~\ref{ch:results}). Within a cell the only remaining source of variation is the decoding randomness $\varepsilon$ of Equation~\eqref{eq:noise}. Each repetition is an independent draw from the same distribution.

\section{Knowledge Base and Query Set}
\label{sec:corpus}

Following Section~\ref{sec:knowledgebase}, both the knowledge base and the query set come from MS MARCO \citep{nguyen2016msmarco}, read from the Hugging Face export of version 2.1. It holds 101,093 queries, of which 12,467 carry a non-empty \texttt{wellFormedAnswers} field. Only those rows are used, because the well-formed answers are the references against which correctness is measured.

\paragraph{The knowledge base.} The corpus holds every passage attached to any of the 12,467 rows. MS MARCO provides at most ten candidate passages per query, retrieved by Bing and then judged by annotators. The corpus therefore contains the gold evidence for every query. Each passage forms one chunk and is used unchanged. 

\paragraph{The query set.} A total of $N_Q = 100$ queries are sampled from the same 12,467 rows. They fall into two categories. A row is \textit{single-passage} when exactly one of its passages is marked \texttt{is\_selected}, and \textit{multi-passage} when two or more are. Fifty queries are drawn from each category.

Within each category the draw is balanced across MS MARCO's five intent types, \textsc{description}, \textsc{entity}, \textsc{location}, \textsc{numeric} and \textsc{person}, ten queries per type. The distribution within MS MARCO is not uniform: on the well-formed rows \textsc{description} makes up 44\% and \textsc{person} 7\%. Of the 12,467 well-formed rows, 12,062 are single-passage and only 386 multi-passage. The remaining 19 rows have no passage marked and are not sampled.

\section{Experimental Procedure}
\label{sec:procedure}

\subsection{Stage 1: Retrieval and the Retrieval Gate}

Queries and passages are encoded with \texttt{all-mpnet-base-v2}, using mean pooling over the last layer and L2 normalisation. For normalised vectors, the inner product therefore equals the cosine similarity.



The retrieval stage is also where the evidence is checked. A query passes the \textit{retrieval gate} when every passage its annotation marks as supporting appears among the $K = 10$ returned. Table~\ref{tab:gate} reports the result.
\begin{table}[htbp]
\centering
\caption{The retrieval gate at the depth used, $K = 10$. A query passes when all of its gold passages are among the ten returned.}
\label{tab:gate}
\begin{tabular}{lrrr}
\toprule
$K$ & Passed & single-passage & multi-passage \\
\midrule
$10$ (used) & 88 / 100 & 45 / 50 & 43 / 50  \\
\bottomrule
\end{tabular}
\end{table}

Queries that fail the gate are excluded from the retrieved level. The measures of Section~\ref{sec:metrics} cannot be meaningfully interpreted when the relevant evidence was never supplied. When the context lacks the passage carrying the answer, a correct answer can only have come from the generator's memory. An incorrect answer cannot be told apart from a faithful reading of the other passages.

\subsection{Stage 2: Generation}

For every query, every cell of the design is run $n = 3$ times, except the retrieved-level cells of queries that failed the gate. The cells are visited in a random order per query. Every generation is an independent request. 

The order of the passages within the context is shuffled, because answer accuracy depends on where the relevant passage is located \citep{liu2024lostmiddle}. Shuffling randomises the position of the gold passages across conditions, so that position cannot systematically favour any condition. 

\section{Evaluation}
\label{sec:metrics}

The evaluation scores each answer on the four properties of Section~\ref{sec:goodanswer}.

Each property is estimated by one measure, except correctness, which is estimated by three. The first three are the measures that the survey of \citet{yu2024ragevalsurvey} names as the evaluation targets of the generation component: faithfulness relates the response to the retrieved documents, answer relevance relates it to the query, and correctness relates it to the reference answer. Retrieval measures are excluded by construction, because the system prompt cannot change what is retrieved.

\subsection{Six Measures for Four Properties}

\paragraph{Faithfulness.} Whether every claim in the answer can be traced to a supplied passage. It is scored reference-free, after \citet{es2023ragas}, in two steps. First the judge (Section~\ref{sec:judge}) breaks the answer into atomic factual claims. Then, for each claim, the judge decides whether the passages state it explicitly.  

\paragraph{Answer relevance.} Whether the answer addresses the question actually asked. A prompt can degrade an answer without making it false. Relevance is scored reference-free, following \citet{es2023ragas}: the judge generates three questions to which the answer would be a suitable response. The score is the mean cosine similarity between these questions and the original question, computed in the embedding space of the retrieval encoder. 

\paragraph{Correctness.} Whether the answer agrees with the reference answers $A_q$. It is measured with the two metrics that MS MARCO uses for its free-form answers \citep{nguyen2016msmarco} and that \citet{lewis2020rag} report on this task: ROUGE-L \citep{lin2004rouge} and BLEU-1 \citep{papineni2002bleu}. Both compare the answer with the references word by word. Before scoring, the text is lowercased and punctuation is replaced by spaces. Citation markers such as \texttt{[3]} are also removed from the answer. Baseline $B$ asks for them, and they would otherwise penalise its answers for following the prompt. ROUGE-L is the F-measure based on the longest common subsequence of the answer and the reference. BLEU-1 is the share of the answer's words that also occur in a reference. Each word is counted at most as often as it appears in any single reference. The result is multiplied by a brevity penalty when the answer is shorter than the reference closest to it in length. The third metric, BERTScore \citep{zhang2020bertscore}, embeds the answer and the reference with a pretrained encoder, matches every word of one to the most similar word of the other by the cosine of their embeddings, and combines the two directions into an F1 score. It is computed with the standard settings for English \citep{liu2019roberta}.

\paragraph{Instruction adherence.} Whether the answer complies with the instructions in the prompt. It is scored as strict accuracy in the sense of \citet{zhou2023ifeval}: an answer is adherent only if it satisfies every checked requirement. The score is binary. 
\subsection{The Judge}
\label{sec:judge}

Faithfulness and answer relevance have no reference string. Therefore, they are scored by a second language model, as \citet{es2023ragas} and \citet{saadfalcon2024ares} do. The judge is \texttt{gpt-4.1-mini}, a different model from the generator, so that the generator never grades its own output. It is run at temperature 0 so that the verdict depends as far as possible on the answer and the passages and not on a random draw. \citet{zheng2023judge} document several biases of language-model judges, among them a preference for longer answers.

\subsection{Estimating the Effect of a Rewording}
\label{sec:analysis}

The quantity estimated is the average treatment effect of Equation~\eqref{eq:estimand}. Here, $M$ is instantiated by each of the six measures in turn. Every measure is oriented so that a higher value is better. A positive effect means the rewording improved that measure.

For one measure $M$, one evidence level $e$ and one baseline $b$, the estimate of the effect of wording $j$ is
\begin{equation}
\label{eq:tauhat}
\hat{\tau}_{b,j} \;=\; \frac{1}{|I|} \sum_{i \in I} \Bigl( \bar{M}_{ij} - \bar{M}_{i0} \Bigr), \qquad \bar{M}_{ij} \;=\; \frac{1}{n} \sum_{r=1}^{n} M\bigl(y_{ijer}\bigr),
\end{equation}
where $I$ is the set of queries and $\bar{M}_{ij}$ is the mean of the measure over the $n$ repetitions of query $i$ under condition $S_{b,j}$. Equation~\eqref{eq:tauhat} is the sample version of Equation~\eqref{eq:estimand}. The inner average estimates the expectation over decoding, and the outer one the expectation over queries. The repetitions are averaged first, so that each query contributes one difference. 


\chapter{Results and Discussion}
\label{ch:results}



The run produced 6,768 answers. At the gold level all 100 queries are usable, which gives $100 \times 12 \times 3 = 3{,}600$ answers. At the retrieved level 88 of the 100 queries pass the gate of Section~\ref{sec:procedure} at $K = 10$, which gives $88 \times 12 \times 3 = 3{,}168$. 
\section{Descriptive Results}
\label{sec:descriptive}

Table~\ref{tab:cellmeans} gives the mean of each measure in each cell, and Table~\ref{tab:length} the mean length of the answers.

\begin{table}[htbp]
\centering
\caption{Mean of each measure per condition. The gold level rests on 100 queries and the retrieved level on the 88 that pass the gate at $K = 10$. Correctness is given as ROUGE-L, BLEU-1 and BERTScore. Adherence is undefined for baseline $A$, whose only requirement is grounding.}
\label{tab:cellmeans}
\setlength{\tabcolsep}{4pt}
\begin{tabular}{lrrrrrr}
\toprule
 & Faithful. & Relevance & ROUGE-L & BLEU-1 & BERTScore & Adherence \\
\midrule
\multicolumn{7}{l}{\textit{gold, $|I| = 100$}} \\
A0 & 0.917 & 0.736 & 0.541 & 0.494 & 0.597 & n/a \\
A1 & 0.922 & 0.730 & 0.519 & 0.463 & 0.583 & n/a \\
A2 & 0.930 & 0.734 & 0.550 & 0.505 & 0.610 & n/a \\
A3 & 0.930 & 0.741 & 0.539 & 0.492 & 0.603 & n/a \\
B0 & 0.912 & 0.696 & 0.362 & 0.285 & 0.440 & 0.223 \\
B1 & 0.942 & 0.705 & 0.407 & 0.332 & 0.470 & 0.447 \\
B2 & 0.924 & 0.703 & 0.397 & 0.319 & 0.464 & 0.343 \\
B3 & 0.931 & 0.703 & 0.403 & 0.327 & 0.473 & 0.297 \\
C0 & 0.924 & 0.729 & 0.262 & 0.164 & 0.358 & 0.943 \\
C1 & 0.919 & 0.738 & 0.272 & 0.174 & 0.359 & 0.940 \\
C2 & 0.941 & 0.728 & 0.296 & 0.193 & 0.387 & 0.853 \\
C3 & 0.939 & 0.731 & 0.322 & 0.214 & 0.395 & 0.923 \\
\addlinespace
\multicolumn{7}{l}{\textit{retrieved, $K = 10$, $|I| = 88$}} \\
A0 & 0.968 & 0.737 & 0.436 & 0.377 & 0.497 & n/a \\
A1 & 0.976 & 0.734 & 0.418 & 0.359 & 0.481 & n/a \\
A2 & 0.974 & 0.732 & 0.443 & 0.383 & 0.508 & n/a \\
A3 & 0.976 & 0.736 & 0.438 & 0.376 & 0.499 & n/a \\
B0 & 0.962 & 0.693 & 0.233 & 0.157 & 0.325 & 0.083 \\
B1 & 0.971 & 0.705 & 0.297 & 0.220 & 0.369 & 0.345 \\
B2 & 0.957 & 0.696 & 0.267 & 0.192 & 0.349 & 0.193 \\
B3 & 0.965 & 0.702 & 0.282 & 0.205 & 0.358 & 0.205 \\
C0 & 0.963 & 0.734 & 0.201 & 0.119 & 0.276 & 0.913 \\
C1 & 0.973 & 0.738 & 0.207 & 0.125 & 0.279 & 0.932 \\
C2 & 0.973 & 0.746 & 0.229 & 0.141 & 0.306 & 0.788 \\
C3 & 0.959 & 0.730 & 0.267 & 0.170 & 0.337 & 0.920 \\
\bottomrule
\end{tabular}
\end{table}

\begin{table}[htbp]
\centering
\caption{Mean answer length in tokens per condition.}
\label{tab:length}
\setlength{\tabcolsep}{3.5pt}
\begin{tabular}{lrrrrrrrrrrrr}
\toprule
 & A0 & A1 & A2 & A3 & B0 & B1 & B2 & B3 & C0 & C1 & C2 & C3 \\
\midrule
gold & 41 & 44 & 38 & 39 & 67 & 59 & 62 & 60 & 96 & 93 & 84 & 68 \\
retrieved & 67 & 80 & 64 & 68 & 112 & 93 & 104 & 99 & 140 & 136 & 119 & 88 \\
\bottomrule
\end{tabular}
\end{table}

The measures differ sharply in how sensitive they are to the prompt. Faithfulness is between $0.912$ and $0.976$ everywhere. Relevance is stable, between $0.693$ and $0.746$. Adherence runs from $0.083$ to $0.943$. ROUGE-L runs from $0.201$ to $0.550$, BLEU-1 from $0.119$ to $0.505$ and BERTScore from $0.276$ to $0.610$. All three are highest under baseline $A$, whose answers are the shortest, and lowest under baseline $C$, whose answers are the longest, because they must include a justification.

Faithfulness is higher at the retrieved level than at the gold level in every condition, although the gold passages are the better evidence.

\section{The Effect of Rewording}
\label{sec:effects}

In the tables that follow, a star marks an estimate whose Holm-adjusted $p$-value is below $0.05$.\footnote{Each estimate $\hat{\tau}_{b,j}$ is tested with a one-sample $t$-test on its paired differences per query. With 102 tests at $\alpha = 0.05$, some would be significant by chance alone, so the $p$-values are adjusted by the procedure of \citet{holm1979}.} The 95\% intervals are not adjusted, so an interval can exclude zero although its estimate carries no star.

\subsection{Instruction Adherence}

The largest effects appear on this measure. Table~\ref{tab:adherence} gives all six estimates at both evidence levels, for baselines $B$ and $C$.
\begin{table}[htbp]
\centering
\caption{Effect of each rewording on instruction adherence, against the canonical form of its own baseline. A star marks a Holm-adjusted $p$-value below $0.05$.}
\label{tab:adherence}
\begin{tabular}{lllll}
\toprule
 & \multicolumn{2}{c}{gold} & \multicolumn{2}{c}{retrieved, $K = 10$} \\
\cmidrule(lr){2-3} \cmidrule(lr){4-5}
Rewording & $\hat{\tau}$ & 95\% CI & $\hat{\tau}$ & 95\% CI \\
\midrule
B1 lexical   & $+0.223$* & $[+0.158, +0.289]$ & $+0.261$* & $[+0.187, +0.336]$ \\
B2 syntactic & $+0.120$* & $[+0.071, +0.169]$ & $+0.110$* & $[+0.045, +0.175]$ \\
B3 format    & $+0.073$* & $[+0.026, +0.121]$ & $+0.121$* & $[+0.058, +0.185]$ \\
C1 lexical   & $-0.003$\phantom{*} & $[-0.039, +0.032]$ & $+0.019$\phantom{*} & $[-0.020, +0.058]$ \\
C2 syntactic & $-0.090$* & $[-0.140, -0.040]$ & $-0.125$* & $[-0.187, -0.063]$ \\
C3 format    & $-0.020$\phantom{*} & $[-0.067, +0.027]$ & $+0.008$\phantom{*} & $[-0.048, +0.063]$ \\
\bottomrule
\end{tabular}
\end{table}

Eight of the twelve estimates are significant. Every rewording of baseline $B$ raises adherence at both evidence levels. Of the rewordings of $C$, only the syntactic one, C2, has a significant effect. C1 and C3 have no significant effect.

Regarding $B$ the rule fails an answer when any of its sentences carries no citation. Under the canonical wording B0 the first sentence of an answer is uncited in 77\% of the gold answers and in 90\% of the retrieved ones. Under B1 these shares fall to 53\% and 63\%. 

For baseline $C$, the rewording that lowers adherence is the one that turns the instructions into plain passive statements, such as \textit{The answer is then given, with an explanation of why it is chosen}. A plausible explanation is that it reads as a description rather than a request, and is therefore followed less often. This fits the concern raised in Section~\ref{sec:conditions}, that such statements may state a requirement less strongly than the canonical imperative.

\subsection{Faithfulness}

One of the nine estimates at the gold level is significant, namely B1, $\hat{\tau} = +0.031$ with a 95\% interval of $[+0.014, +0.047]$. The same comparison at the retrieved level gives $+0.009$ with an interval that contains zero. A result that appears at one level and not the other is weak evidence.

\subsection{Answer Relevance}

No estimate is significant at either level. The largest estimate in absolute value is $+0.012$.

\subsection{Correctness}

Tables~\ref{tab:rouge}, \ref{tab:bleu} and~\ref{tab:bertscore} give the estimates.

\begin{table}[htbp]
\centering
\caption{Effect of each rewording on correctness measured by ROUGE-L, against the canonical form of its own baseline. A star marks a Holm-adjusted $p$-value below $0.05$.}
\label{tab:rouge}
\begin{tabular}{lllll}
\toprule
 & \multicolumn{2}{c}{gold} & \multicolumn{2}{c}{retrieved, $K = 10$} \\
\cmidrule(lr){2-3} \cmidrule(lr){4-5}
Rewording & $\hat{\tau}$ & 95\% CI & $\hat{\tau}$ & 95\% CI \\
\midrule
A1 lexical   & $-0.022$* & $[-0.036, -0.008]$ & $-0.018$* & $[-0.031, -0.005]$ \\
A2 syntactic & $+0.009$\phantom{*} & $[-0.002, +0.021]$ & $+0.008$\phantom{*} & $[-0.006, +0.022]$ \\
A3 format    & $-0.002$\phantom{*} & $[-0.016, +0.012]$ & $+0.002$\phantom{*} & $[-0.007, +0.011]$ \\
B1 lexical   & $+0.045$* & $[+0.027, +0.062]$ & $+0.064$* & $[+0.042, +0.087]$ \\
B2 syntactic & $+0.034$* & $[+0.020, +0.049]$ & $+0.034$* & $[+0.016, +0.051]$ \\
B3 format    & $+0.041$* & $[+0.027, +0.054]$ & $+0.049$* & $[+0.028, +0.070]$ \\
C1 lexical   & $+0.010$\phantom{*} & $[-0.000, +0.020]$ & $+0.006$\phantom{*} & $[-0.002, +0.013]$ \\
C2 syntactic & $+0.034$* & $[+0.025, +0.043]$ & $+0.027$* & $[+0.020, +0.034]$ \\
C3 format    & $+0.060$* & $[+0.047, +0.072]$ & $+0.066$* & $[+0.054, +0.077]$ \\
\bottomrule
\end{tabular}
\end{table}

\begin{table}[htbp]
\centering
\caption{Effect of each rewording on correctness measured by BLEU-1, against the canonical form of its own baseline. A star marks a Holm-adjusted $p$-value below $0.05$.}
\label{tab:bleu}
\begin{tabular}{lllll}
\toprule
 & \multicolumn{2}{c}{gold} & \multicolumn{2}{c}{retrieved, $K = 10$} \\
\cmidrule(lr){2-3} \cmidrule(lr){4-5}
Rewording & $\hat{\tau}$ & 95\% CI & $\hat{\tau}$ & 95\% CI \\
\midrule
A1 lexical   & $-0.031$* & $[-0.049, -0.013]$ & $-0.017$\phantom{*} & $[-0.031, -0.003]$ \\
A2 syntactic & $+0.011$\phantom{*} & $[-0.003, +0.025]$ & $+0.007$\phantom{*} & $[-0.009, +0.022]$ \\
A3 format    & $-0.002$\phantom{*} & $[-0.020, +0.015]$ & $-0.001$\phantom{*} & $[-0.011, +0.009]$ \\
B1 lexical   & $+0.047$* & $[+0.027, +0.068]$ & $+0.063$* & $[+0.037, +0.089]$ \\
B2 syntactic & $+0.034$* & $[+0.018, +0.050]$ & $+0.035$* & $[+0.016, +0.054]$ \\
B3 format    & $+0.042$* & $[+0.026, +0.058]$ & $+0.048$* & $[+0.028, +0.069]$ \\
C1 lexical   & $+0.010$\phantom{*} & $[+0.002, +0.018]$ & $+0.006$\phantom{*} & $[+0.001, +0.010]$ \\
C2 syntactic & $+0.029$* & $[+0.022, +0.037]$ & $+0.022$* & $[+0.017, +0.027]$ \\
C3 format    & $+0.050$* & $[+0.041, +0.059]$ & $+0.051$* & $[+0.044, +0.058]$ \\
\bottomrule
\end{tabular}
\end{table}

\begin{table}[htbp]
\centering
\caption{Effect of each rewording on correctness measured by BERTScore, against the canonical form of its own baseline. A star marks a Holm-adjusted $p$-value below $0.05$.}
\label{tab:bertscore}
\begin{tabular}{lllll}
\toprule
 & \multicolumn{2}{c}{gold} & \multicolumn{2}{c}{retrieved, $K = 10$} \\
\cmidrule(lr){2-3} \cmidrule(lr){4-5}
Rewording & $\hat{\tau}$ & 95\% CI & $\hat{\tau}$ & 95\% CI \\
\midrule
A1 lexical   & $-0.014$* & $[-0.024, -0.004]$ & $-0.016$* & $[-0.027, -0.005]$ \\
A2 syntactic & $+0.013$* & $[+0.004, +0.022]$ & $+0.011$\phantom{*} & $[+0.000, +0.021]$ \\
A3 format    & $+0.006$\phantom{*} & $[-0.004, +0.015]$ & $+0.002$\phantom{*} & $[-0.007, +0.011]$ \\
B1 lexical   & $+0.030$* & $[+0.017, +0.044]$ & $+0.044$* & $[+0.027, +0.062]$ \\
B2 syntactic & $+0.024$* & $[+0.013, +0.035]$ & $+0.024$* & $[+0.011, +0.037]$ \\
B3 format    & $+0.033$* & $[+0.022, +0.044]$ & $+0.033$* & $[+0.018, +0.048]$ \\
C1 lexical   & $+0.002$\phantom{*} & $[-0.006, +0.009]$ & $+0.003$\phantom{*} & $[-0.005, +0.011]$ \\
C2 syntactic & $+0.030$* & $[+0.021, +0.038]$ & $+0.030$* & $[+0.021, +0.038]$ \\
C3 format    & $+0.037$* & $[+0.025, +0.049]$ & $+0.061$* & $[+0.047, +0.076]$ \\
\bottomrule
\end{tabular}
\end{table}

Every rewording of baseline $B$ raises correctness. Of the rewordings of $C$, the syntactic and format ones raise it. The format one raises it most. No estimate for the lexical rewording C1 is significant. For baseline $A$, the lexical rewording A1 lowers it. The BLEU-1 estimate there is not significant after adjustment. Every significant BERTScore effect has the same direction as the ROUGE-L and BLEU-1 effects of that rewording. 

These effects track answer length. Every rewording that raised correctness produced shorter answers than its canonical form (Table~\ref{tab:length}), and A1, which lowered it, produced longer ones. A2, which raised BERTScore, produced slightly shorter answers as well. Across the eighteen comparisons, the change in mean length and the change in ROUGE-L have a correlation of $-0.79$, and for BERTScore the correlation is $-0.87$. The reason is that the references are short, so a shorter answer shares a larger fraction of its words, and of its meaning, with them. BERTScore is also affected, because its precision component penalises an answer for everything it says beyond the reference. The recall component of ROUGE-L alone, which measures how much of the reference an answer contains and does not penalise length, shows no significant gain for any rewording. The rewordings therefore changed how closely the answers matched the reference mainly by changing their length, not their content.

\section{The Verdict on the Hypothesis}
\label{sec:verdict}

One hundred and two estimates were made across the six measures of Section~\ref{sec:metrics}, and forty-five of them are significant after the Holm adjustment: eight on adherence, thirty-six on the three correctness metrics and one on faithfulness. $H_0$ is therefore rejected. $H_1$, that at least one rewording changes answer quality, is supported: seven of the nine rewordings have a significant effect on at least one measure. A3 and C1 have no significant effect on any measure.

The rejection does not rest on the single faithfulness result. It rests on instruction adherence, where four rewordings have clear effects, three positive and one negative, and on correctness, where the effects are small and track answer length.

\section{Limitations}
\label{sec:limits}
Four limitations apply to these results. The first is the collection. Every result is conditional on MS MARCO. The effects therefore suggest, but do not establish, what may be expected on other data. The second concerns correctness. ROUGE-L and BLEU-1 reward word overlap and BERTScore semantic similarity. None of them measures factual accuracy, and all three favour answers close in length to the reference. The third concerns the rewordings themselves. Four places where a rewording may have changed how strongly a requirement is stated were named in Section~\ref{sec:conditions}. The effect of C2 in particular may partly reflect a weaker statement of the requirement rather than a change of wording alone. The fourth concerns the judge: faithfulness and answer relevance are scored by a language model, which may be biased, for example towards longer answers (Section~\ref{sec:judge}).


\chapter{Conclusion and Future Work}
\label{ch:conclusion}

\section{Conclusion}

This thesis asked to what extent systematic changes in the system prompt influence the behaviour and outcomes of a RAG system when the user input, the knowledge base, the retrieval mechanism and the underlying language model are held fixed. It answered the question for one kind of change, a change of wording that leaves the requirements of the prompt unchanged.

The answer is that wording matters. It mattered most for whether the answers complied with the prompt. Rewording the citation instruction of \citet{gao2023alce} raised the share of fully cited answers. Turning the instruction of \citet{mao2024fitrag} into passive statements lowered the share of justified answers. Wording also changed the length of the answers, and thereby their agreement with the reference answer. It had no consistent effect on whether the answers stayed within the supplied passages or addressed the question. These results held at both evidence levels, that is, both with only the supporting passages and with the actual output of the retriever.

The contribution of this thesis: The formal model of Chapter 4 makes the system prompt an explicit variable of the generator and defines its effect as an average treatment effect over queries. The design of Chapter 5 separates what a prompt asks for from how it is worded, fixes everything else, sets the evidence directly, and pairs every comparison within a query. The same design can be reused with other prompts, models and datasets.

\section{Future Work}
 

\paragraph{More generators.} Every result here comes from one model, \texttt{gpt-4o-mini}. Whether other models respond to the same rewordings in the same way remains an open question. Repeating the experiment with a second API model and an open model of a different size would show whether the effects are specific to the model or arise from the wording itself.

\paragraph{Correctness beyond overlap.} The three correctness metrics compare wording and meaning with a reference but do not verify factual accuracy. A judge that decides whether an answer contains the fact stated in the reference, validated against human labels on a sample, would measure correctness independently of answer length.

\paragraph{Other data.}  A dataset with longer answers, and an open collection in which the answer may be missing, would test whether the effects hold where retrieval is harder and where the prompt has to govern what the model does when the evidence does not suffice.

\bibliographystyle{apalike}
\bibliography{thesis}

\end{document}

